\documentclass[11pt,a4paper]{article}

\usepackage[margin=1in]{geometry}
\usepackage{amsmath,amssymb}
\usepackage{booktabs}
\usepackage{graphicx}
\usepackage{hyperref}
\usepackage{xcolor}
\usepackage{enumitem}
\usepackage{listings}
\usepackage[T1]{fontenc}
\usepackage{lmodern}
\usepackage{microtype}

\hypersetup{colorlinks=true,linkcolor=blue!60!black,citecolor=green!50!black,urlcolor=blue!70!black}

\lstset{
  basicstyle=\ttfamily\small,
  breaklines=true,
  frame=single,
  backgroundcolor=\color{gray!8},
  keywordstyle=\color{blue!70!black},
  commentstyle=\color{green!50!black},
}

\title{%
  Technical Notes: Depth-Based Student Distillation\\
  for Quadruped Parkour on Genesis\\[6pt]
  \large LeggedGym-Ex Repository Status Report
}
\author{Auto-generated Research Note}
\date{March 25, 2026}

\begin{document}
\maketitle

% ============================================================================
\section{Introduction and Goals}

This document records the current state of the \texttt{LeggedGym-Ex} repository,
which implements a two-phase teacher--student training pipeline for quadruped
parkour locomotion on the Unitree Go2 robot, using the Genesis physics simulator.

\paragraph{End goal.}
Train a student policy that receives \emph{monocular RGB} images, runs them through
a depth-estimation network (e.g.\ Depth Anything V2), and feeds the inferred
depth map to the policy---enabling deployment without a physical depth camera.

\paragraph{Current milestone.}
The student policy is being trained on \emph{ground-truth simulator depth} to
establish a convergent distillation baseline before introducing the
depth-estimation network.

\paragraph{Reference work.}
The architecture and training procedure follow Cheng et al.\ ``Extreme Parkour
with Legged Robots'' (arXiv:2309.14341)~\cite{cheng2023extreme}, with
adaptations for the Genesis simulator and Go2 robot morphology.

% ============================================================================
\section{System Overview}

\subsection{Hardware and Simulator}

\begin{itemize}[nosep]
  \item \textbf{GPU:} 4$\times$ NVIDIA L40S (45\,GB each).
  \item \textbf{Simulator:} Genesis (GPU-accelerated rigid-body physics with
        batch depth rendering via CUDA scene cameras on Linux).
  \item \textbf{Robot:} Unitree Go2 (12 DoF, 4 legs $\times$ 3 joints).
\end{itemize}

\subsection{Two-Phase Training Pipeline}

\begin{enumerate}[nosep]
  \item \textbf{Phase 1 --- Teacher (PPO):}
        A privileged teacher policy observes proprioception \emph{plus} 132
        ray-cast scandot heights and is trained with PPO on a parkour curriculum
        (stairs, hurdle blocks, gaps). The critic additionally observes base
        linear velocity, dynamics parameters, and contact states.
  \item \textbf{Phase 2 --- Student (DAgger distillation):}
        A student policy observes only proprioception (53-dim) and a single
        $58 \times 87$ depth image. A CNN--GRU depth encoder produces a
        32-dim terrain latent and a 2-dim yaw estimate. The student's actor MLP
        is initialized from the teacher's actor weights and fine-tuned via
        DAgger: at each step the frozen teacher provides target actions, and the
        student minimizes L2 action loss + L2 yaw loss.
\end{enumerate}

% ============================================================================
\section{Architecture Details}

\subsection{Observation Layout}

The proprioceptive observation vector (53-dim) is ordered as:

\begin{center}
\begin{tabular}{llr}
\toprule
Component & Description & Dims \\
\midrule
Commands & goal vel, heading ($\sin\theta_e$, $\cos\theta_e$), goal speed & 7 \\
Gravity & projected gravity in body frame & 3 \\
Angular velocity & base angular velocity (scaled) & 3 \\
Joint positions & $q - q_{\mathrm{default}}$ (scaled) & 12 \\
Joint velocities & $\dot{q}$ (scaled) & 12 \\
Previous actions & $a_{t-1}$ & 12 \\
Foot contacts & binary contact indicators & 4 \\
\bottomrule
\end{tabular}
\end{center}

The teacher additionally receives 132 scandot heights (12 $\times$ 11 grid)
appended after the proprioceptive vector, for a total actor observation of 185
dims. The student's actor only sees the 53-dim proprio; terrain perception
comes exclusively through the depth encoder's 32-dim latent.

Heading command indices 4--5 ($\sin\theta_e$, $\cos\theta_e$) are
\emph{masked to zero} in the depth encoder's proprioceptive input, forcing the
GRU to predict yaw from the depth stream.

\subsection{Student Network Architecture}

\begin{center}
\begin{tabular}{lll}
\toprule
Module & Architecture & Output \\
\midrule
\texttt{DepthBackbone} & Conv2d(1,32,k=5) $\to$ MaxPool(2) $\to$ ELU
  $\to$ Conv2d(32,64,k=3) $\to$ ELU $\to$ FC(128) $\to$ FC(32) & 32-dim \\
\texttt{CombinationMLP} & FC(32+53, 128) $\to$ ELU $\to$ FC(128, 32) & 32-dim \\
\texttt{GRU} & GRU(input=32, hidden=512, 1 layer) & 512-dim \\
\texttt{OutputMLP} & FC(512, 34) $\to$ Tanh & 32+2 dim \\
\texttt{Actor} & FC(53+32, 512) $\to$ ELU $\to$ FC(256) $\to$ ELU
  $\to$ FC(128) $\to$ ELU $\to$ FC(12) $\to$ Hardtanh & 12-dim \\
\bottomrule
\end{tabular}
\end{center}

The Tanh output is split into a 32-dim \emph{latent} (fed to the actor
alongside proprio) and a 2-dim \emph{raw yaw} estimate, scaled by 1.5$\times$
to match oracle heading magnitude.

\subsection{Mixture of Teacher and Student (MTS)}

Following the reference, a soft gating mechanism replaces the student's heading
observation with its own yaw prediction only when the prediction is close to
the oracle:
\[
  \mathbf{o}_{\mathrm{heading}} =
  \begin{cases}
    \hat{\mathbf{y}}_{\mathrm{student}} & \text{if } \|\hat{\mathbf{y}}_{\mathrm{student}} - \mathbf{y}_{\mathrm{oracle}}\| < \tau \\
    \mathbf{y}_{\mathrm{oracle}} & \text{otherwise}
  \end{cases}
\]
with threshold $\tau = 0.6$. This prevents early-training distribution shift
while encouraging the student to eventually rely on its own yaw prediction.

% ============================================================================
\section{Differences from Canonical Implementations}

Table~\ref{tab:diffs} summarizes the key differences between our implementation
and the reference repositories.

\begin{table}[ht]
\centering
\caption{Comparison with reference implementations.}
\label{tab:diffs}
\small
\begin{tabular}{p{3.2cm}p{4.2cm}p{4.2cm}}
\toprule
\textbf{Aspect} & \textbf{Extreme Parkour (IsaacGym)} & \textbf{Ours (Genesis)} \\
\midrule
Simulator & IsaacGym (NVIDIA) & Genesis (open-source, CUDA batch cameras) \\
BPTT strategy & Full 120-step, single \texttt{backward()} & Windowed BPTT(24), 5 windows per rollout \\
\texttt{num\_envs} (student) & 192 & 48 (Genesis depth cameras are memory-intensive) \\
Latent loss & Not used & Not used (removed; was present in initial buggy code) \\
Depth noise & Per-pixel Gaussian (\texttt{dis\_noise}) & Scalar uniform noise per frame (simplified) \\
Depth normalization & $[0, 1]$ range & $[-0.5, +0.5]$ range \\
Depth resolution & $58 \times 87$ & $58 \times 87$ (identical after crop/resize) \\
Teacher critic & Same as ours & Proprio + scandots + base lin vel + dynamics + contacts \\
Terrain families & Parkour-specific obstacles & Stairs, hurdle blocks, gaps (3 families $\times$ 4 variants) \\
Curriculum & Row-based difficulty & Row-based difficulty (identical scheme) \\
GRU reset on done & No reset during training & No reset during training (matching reference) \\
\bottomrule
\end{tabular}
\end{table}

\paragraph{Why windowed BPTT instead of full BPTT.}
Genesis batch depth rendering allocates substantially more GPU memory per
environment than IsaacGym's tensor-based API. With 48 environments, the
simulator alone consumes ${\sim}11$\,GB on an L40S. Full BPTT through 120
CNN+GRU steps requires retaining the entire computation graph, which caused
silent CUDA hangs when attempted (the process continued consuming CPU/GPU
resources but produced no output and never completed the backward pass). A
24-step window (covering ${\sim}5$ depth camera updates at decimation=5)
provides meaningful temporal gradient flow while bounding peak memory.

\paragraph{Why 48 environments.}
The reference uses 192 environments on IsaacGym. Initial attempts with 192
environments on Genesis caused immediate OOM during the first CNN forward pass
(44\,GB GPU fully consumed by simulator state alone). Reducing to 48
environments leaves sufficient headroom for the BPTT computation graph.

% ============================================================================
\section{Bugs Found and Fixed}

The initial codebase had five bugs preventing student convergence, discovered
through careful comparison with the reference implementation.

\subsection{Bug 1 (Critical): Per-Step Backward Destroyed GRU BPTT}

The original \texttt{ParkourDistillation.accumulate\_step()} called
\texttt{backward()} at \emph{each} of the 120 rollout steps, followed by
\texttt{detach\_gru\_hidden()}. This is truncated BPTT(1)---the GRU receives
zero temporal gradient and cannot learn any sequential dependencies.

\textbf{Fix:} Rewrote the distillation algorithm to collect live-graph tensors
across a configurable window (\texttt{bptt\_window=24}), compute loss once per
window, and call \texttt{backward()} at window boundaries. The GRU hidden state
is detached only after each window's backward pass. Gradients accumulate across
all 5 windows, with a single optimizer step after the full 120-step rollout.

\subsection{Bug 2 (High): Teacher Observations Lacked Noise}

The student environment applied observation noise to \texttt{obs\_buf} (the
student's proprioceptive input) but constructed \texttt{teacher\_actor\_obs\_buf}
from clean sensor readings. Since the teacher was \emph{trained with noise},
querying it with clean observations produced misaligned target actions.

\textbf{Fix:} Applied the same noise realization to both the student's
proprioceptive observation and the proprioceptive portion of the teacher's
observation. Scandot noise (matching the teacher's training noise scale) is also
applied to the teacher's scandot portion.

\subsection{Bug 3 (Medium): Latent Loss with Range Mismatch}

The distillation loss included an L2 term between the student's depth latent
(Tanh, range $[-1,1]$) and the teacher's scandot latent (ELU, range
$(-1,+\infty)$). The reference uses \emph{only} action loss + yaw loss. The
range mismatch meant the student could never perfectly match teacher latents,
creating a persistent conflicting gradient.

\textbf{Fix:} Removed the latent loss entirely. The student is trained only on
action matching and yaw matching, consistent with the reference.

\subsection{Bug 4 (Medium): Excessive Environment Count}

The student config specified \texttt{num\_envs=2048}, inherited from the teacher.
With full BPTT (or even windowed BPTT), this caused OOM since Genesis depth
cameras consume significant GPU memory per environment.

\textbf{Fix:} Reduced to 48 environments (the maximum that fits on one L40S
with comfortable headroom for the backward pass).

\subsection{Bug 5 (Low): Missing Depth Noise}

The reference applies Gaussian noise (\texttt{dis\_noise}) to depth images for
sim-to-real robustness. The initial codebase had no depth augmentation.

\textbf{Fix:} Added configurable depth noise (\texttt{depth\_noise\_level=0.1})
applied as a scalar uniform perturbation after sensor rendering.

\subsection{Additional Fix: In-Place Depth Buffer Modification}

During full BPTT development, the backward pass failed with
``\texttt{RuntimeError: variable has been modified by an inplace operation}''
because the simulator's depth buffer is modified in-place between steps by
\texttt{update\_sensors()}. The student's \texttt{forward\_depth()} received a
reference to the same buffer at each step, but autograd expected the original
tensor.

\textbf{Fix:} Clone the depth tensor before each forward pass:
\texttt{student\_depth.clone()}.

% ============================================================================
\section{Training Results}

\subsection{Teacher Training}

The teacher was trained for 2500 iterations (4096 envs, PPO, ${\sim}2$ hours on
one L40S). Table~\ref{tab:teacher} shows the convergence trajectory.

\begin{table}[ht]
\centering
\caption{Teacher training metrics (PPO, 4096 envs).}
\label{tab:teacher}
\begin{tabular}{rcccc}
\toprule
Iter & Reward & Success & Progress & Terrain Level \\
\midrule
0    & $-1.1$ & 0.000 & 0.008 & 0.00 \\
500  & 51.4   & 0.827 & 0.907 & 2.03 \\
1000 & 52.9   & 0.791 & 0.908 & 1.99 \\
1500 & 57.5   & 0.685 & 0.937 & 2.03 \\
2000 & 74.4   & 0.225 & 0.926 & 1.99 \\
\bottomrule
\end{tabular}
\end{table}

The teacher reached 83\% success rate at peak (iter 500), with mean reward
climbing to 74 as the curriculum pushed harder terrain. Progress ratio stabilized
above 90\%.

\subsection{Student Training --- Run 1 (Iterations 0--1003)}

The student was trained with windowed BPTT(24), 48 envs, 120 steps/env,
lr=$10^{-3}$, on one L40S GPU. Table~\ref{tab:student1} shows the convergence.

\begin{table}[ht]
\centering
\caption{Student Run 1 metrics (DAgger distillation, 48 envs, windowed BPTT).}
\label{tab:student1}
\small
\begin{tabular}{rccccccccc}
\toprule
Iter & Act.\ Loss & Yaw & Reward & Success & Progress & Level & Row & WPs & Ep.\ Len \\
\midrule
0    & 1.562 & 0.284 &   4.3 & 0.000 & 0.062 & 0.00 & 0.00 & 0.0 &  50 \\
100  & 0.748 & 0.108 &  46.9 & 0.092 & 0.710 & 1.00 & 0.98 & 1.5 & 549 \\
200  & 0.477 & 0.065 &  72.6 & 0.000 & 0.627 & 1.47 & 1.43 & 1.8 & 773 \\
300  & 0.425 & 0.052 &  82.7 & \textbf{0.471} & 0.939 & 2.04 & 1.91 & 3.0 & 861 \\
500  & 0.403 & 0.041 &  86.6 & 0.025 & 0.811 & 2.00 & 1.65 & 1.7 & 901 \\
700  & 0.313 & 0.038 &  88.9 & 0.142 & 0.925 & 1.73 & 2.06 & 2.8 & 928 \\
800  & 0.324 & 0.031 &  \textbf{92.7} & 0.100 & \textbf{0.949} & 2.09 & 1.28 & 2.2 & 959 \\
1000 & 0.329 & 0.041 &  89.5 & 0.358 & 0.946 & 1.91 & 2.15 & 3.0 & 928 \\
\bottomrule
\end{tabular}
\end{table}

\paragraph{Key observations.}
\begin{itemize}[nosep]
  \item Action loss converged from 1.56 to ${\sim}0.33$ (5$\times$ reduction).
  \item Yaw loss converged from 0.28 to ${\sim}0.03$ (10$\times$ reduction).
  \item Mean reward reached 92.7 at iter 800, \emph{exceeding} the teacher's
        typical 74 reward---likely because the student's depth perception
        provides smoother terrain information than the teacher's sparse scandots.
  \item Peak success rate of 47\% (iter 300) and 80\% (measured via TensorBoard
        max) with progress ratio consistently above 94\%.
  \item Terrain curriculum advanced to level ${\sim}2.0$ (out of 4 rows), with
        the student completing nearly full courses on medium-difficulty terrain.
\end{itemize}

\subsection{Student Training --- Run 2 (Resumed, Iterations 1003--1492)}

After resuming from the Run~1 checkpoint, the student continued improving.
Table~\ref{tab:student2} shows metrics relative to the resumed run's local
iteration counter.

\begin{table}[ht]
\centering
\caption{Student Run 2 (resumed from iter 1003).}
\label{tab:student2}
\begin{tabular}{rccccc}
\toprule
$\Delta$Iter & Act.\ Loss & Yaw & Reward & Success & Progress \\
\midrule
0     & 0.346 & 0.024 &  2.7  & 0.000 & 0.057 \\
100   & 0.314 & 0.028 & 92.1  & 0.158 & 0.954 \\
200   & 0.286 & 0.025 & 92.5  & 0.017 & 0.974 \\
300   & 0.292 & 0.030 & 93.0  & 0.067 & 0.944 \\
400   & 0.316 & 0.027 & 93.5  & 0.083 & 0.939 \\
489   & \textbf{0.272} & \textbf{0.017} & 92.8  & 0.000 & 0.918 \\
\bottomrule
\end{tabular}
\end{table}

Action loss continued decreasing to 0.27, yaw loss to 0.017, and reward
stabilized in the 92--94 range. The terrain level advanced to ${\sim}2.2$.
The training is still running as of this writing.

% ============================================================================
\section{Open Issues and Next Steps}

\subsection{Remaining Issues}

\begin{enumerate}
  \item \textbf{Terminal output buffering.}
        Console output (stdout) stops being captured after ${\sim}500$ iterations,
        though the training process continues running normally. TensorBoard
        events and checkpoints are written correctly. This appears to be a
        Python stdout buffering issue when output is redirected to a pipe.
        Workaround: monitor via TensorBoard or periodic checkpoint inspection.

  \item \textbf{Depth noise is scalar, not per-pixel.}
        The current depth noise augmentation applies a single random scalar
        offset to the entire depth image, rather than per-pixel Gaussian noise
        as in the reference. This provides less diversity for sim-to-real
        robustness. The per-pixel version was deferred to avoid additional GPU
        memory overhead during BPTT.

  \item \textbf{Low environment count (48 vs.\ 192).}
        Genesis's memory-intensive depth rendering limits us to 48 parallel
        environments, compared to 192 in the reference. This reduces sample
        diversity per iteration and may slow convergence. Potential mitigations:
        (a) gradient checkpointing through the CNN to reduce activation memory,
        (b) using Genesis's lower-level API to share rendering resources,
        (c) mixed-precision training (FP16 activations).

  \item \textbf{Success rate instability.}
        Success rate oscillates between 0\% and 47\% across iterations due to
        the small number of episodes completed per iteration (48 envs $\times$
        ${\sim}0.13$ episodes/env/iter $\approx$ 6 episodes). With only 6
        Bernoulli samples, the variance is inherently high. The progress ratio
        (0.94--0.97) is a more stable convergence indicator.

  \item \textbf{Resume path resolution.}
        The \texttt{--resume} flag resolves \texttt{load\_run=-1} by sorting all
        directories in the log root alphabetically and selecting the last one. A
        \texttt{datasets/} directory (used for demo replay) sorts after the
        timestamped run directories (e.g.\ \texttt{Mar24\_...}), causing resume
        to fail. Fix: exclude non-run directories or specify \texttt{--load\_run}
        explicitly.
\end{enumerate}

\subsection{Planned Next Steps}

\begin{enumerate}
  \item \textbf{Complete 5000-iteration training} and evaluate final success rate
        and progress ratio on the hardest terrain rows.

  \item \textbf{Per-pixel depth noise.} Replace scalar noise with per-pixel
        Gaussian noise ($\sigma = 0.1$) matching the reference's
        \texttt{dis\_noise} implementation.

  \item \textbf{Depth estimation network integration.}
        Replace simulator depth with inferred depth from Depth Anything V2
        (or similar monocular depth estimator) applied to the simulator's RGB
        output. This requires:
        \begin{itemize}[nosep]
          \item Adding an RGB camera to the student environment (same pose as
                depth camera).
          \item Running the depth estimator at the same decimation rate as the
                depth camera.
          \item Handling scale/shift ambiguity between metric depth and
                estimator output (affine alignment or learned adapter).
        \end{itemize}

  \item \textbf{Evaluate sim-to-real gap.} Deploy the trained student on
        physical Go2 hardware with an onboard RGB camera and assess transfer
        quality.
\end{enumerate}

% ============================================================================
\section{Hyperparameter Summary}

\begin{table}[ht]
\centering
\caption{Key hyperparameters.}
\begin{tabular}{llcc}
\toprule
& Parameter & Teacher & Student \\
\midrule
\multirow{4}{*}{Env}
& \texttt{num\_envs} & 4096 & 48 \\
& \texttt{episode\_length\_s} & 20 & 20 \\
& \texttt{num\_steps\_per\_env} & 48 & 120 \\
& Depth camera decimation & --- & 5 \\
\midrule
\multirow{4}{*}{Training}
& Algorithm & PPO & DAgger (windowed BPTT) \\
& Learning rate & $10^{-3}$ & $10^{-3}$ \\
& \texttt{max\_grad\_norm} & 1.0 & 1.0 \\
& \texttt{bptt\_window} & --- & 24 \\
\midrule
\multirow{3}{*}{Loss}
& Action loss coef & --- & 1.0 \\
& Yaw loss coef & --- & 1.0 \\
& MTS yaw threshold & --- & 0.6 \\
\midrule
\multirow{2}{*}{Noise}
& Observation noise & Yes & Yes (matched to teacher) \\
& Depth noise level & --- & 0.1 \\
\midrule
\multirow{2}{*}{Terrain}
& Families & \multicolumn{2}{c}{stairs, hurdle\_block, gap} \\
& Rows $\times$ Cols & \multicolumn{2}{c}{$4 \times 12$ (4 difficulty $\times$ 3 families $\times$ 4 variants)} \\
\bottomrule
\end{tabular}
\end{table}

% ============================================================================
\section{Repository Structure (Relevant Files)}

\begin{lstlisting}[language={},frame=none,backgroundcolor=\color{white}]
LeggedGym-Ex/
  legged_gym/
    envs/go2/
      go2_parkour_teacher/
        go2_parkour_teacher.py          # Teacher env (LeggedRobotParkour)
        go2_parkour_teacher_config.py   # Teacher config (PPO, 4096 envs)
      go2_parkour_student/
        go2_parkour_student.py          # Student env (depth + proprio)
        go2_parkour_student_config.py   # Student config (DAgger, 48 envs)
    envs/base/
      legged_robot_parkour.py           # Base parkour env with curriculum
      parkour_observation.py            # Observation spec (prop/actor/critic dims)
    simulator/
      genesis_simulator.py              # Genesis backend (depth rendering)
  rsl_rl/
    algorithms/
      parkour_distillation.py           # Windowed BPTT DAgger algorithm
    runners/
      parkour_student_runner.py         # Student training loop
    modules/
      actor_critic_parkour.py           # Teacher network (scandot encoder)
      actor_critic_parkour_student.py   # Student network (CNN+GRU+actor)
  logs/
    go2_parkour_teacher/                # Teacher checkpoints and TB events
    go2_parkour_student/                # Student checkpoints and TB events
\end{lstlisting}

% ============================================================================
\begin{thebibliography}{1}
\bibitem{cheng2023extreme}
  X.~Cheng, K.~Shi, A.~Agarwal, D.~Pathak,
  ``Extreme Parkour with Legged Robots,''
  \emph{arXiv preprint arXiv:2309.14341}, 2023.
  Code: \url{https://github.com/chengxuxin/extreme-parkour}
\end{thebibliography}

\end{document}
